\section{Exploratory Data Analysis Plan}

\subsection{How Churn-Rate Charts Should Be Read}
In churn-rate charts, a bar does not represent a distribution that sums to
100\% across groups. Each bar represents the churn rate inside one group:
\[
  \text{Churn Rate}_{group}
  =
  \frac{\text{Number of churned customers in the group}}
       {\text{Total customers in the group}}
  \times 100\%.
\]
This means each group should be compared independently against the baseline
churn rate of approximately 9.92\%.

\subsection{Customer Experience Features}
\begin{longtable}{p{0.28\textwidth}p{0.62\textwidth}}
\caption{Planned EDA interpretation by feature}\\
\toprule
\textbf{Feature} & \textbf{Preliminary interpretation} \\
\midrule
\endfirsthead
\toprule
\textbf{Feature} & \textbf{Preliminary interpretation} \\
\midrule
\endhead
Customer satisfaction &
Higher satisfaction tends to lean toward non-churn. Low satisfaction groups are
expected to have higher churn risk. \\
Number of complaints &
More complaints are associated with a higher churn rate and should be treated as
a strong customer-experience risk signal. \\
Number of service calls &
More service calls may indicate unresolved service problems and higher churn
risk. \\
Late payments &
More late payments are associated with higher churn rate, although the signal is
weaker than satisfaction and complaints. \\
Average monthly GB &
Usage groups have relatively similar churn rates; this feature is weak as a
standalone churn driver. \\
Days since last interaction &
Churn rates are expected to be nearly flat across time-since-interaction
groups, so this feature should be combined with other signals. \\
Credit score &
Credit-score groups show small churn-rate differences and are not expected to be
a major standalone churn signal. \\
\bottomrule
\end{longtable}

\subsection{Planned EDA Visuals}
For each selected feature, the final report should include two complementary
visuals:
\begin{enumerate}
  \item Customer count by churn status for each group.
  \item Churn rate for each group.
\end{enumerate}
This structure separates customer distribution from churn risk. A large segment
does not automatically mean high churn risk, and a small segment can still be
important if its churn rate is much higher than the baseline.

\subsection{Priority of Feature Interpretation}
The feature interpretation should be ordered by actionability. Customer
experience and payment-behavior variables are more useful for retention
planning because they can trigger operational actions. Usage and credit
variables should be treated as supporting context unless they show a clear
above-baseline churn-rate pattern.

\begin{enumerate}
  \item Customer satisfaction.
  \item Number of complaints.
  \item Number of service calls.
  \item Late payments.
  \item Average monthly GB.
  \item Days since last interaction.
  \item Credit score.
\end{enumerate}
